\documentclass[11pt, a4paper]{article}

%\usepackage{geometry}
%\usepackage{epsfig}
%\usepackage{epstopdf}
%\usepackage{graphics}
\usepackage{graphicx}

\setlength{\oddsidemargin}{0cm}
\setlength{\evensidemargin}{0cm}
\setlength{\topmargin}{-1cm}
\setlength{\textheight}{23cm}
\setlength{\textwidth}{16cm}

\pagestyle{headings}

%--------------------------------------------------------

\begin{document}

% --------------------------------------------------------
% title
% --------------------------------------------------------

\thispagestyle{empty}

\begin{figure}
\includegraphics[scale=0.25]{bat.eps}
\end{figure}

\vspace{1.0cm}

\begin{center}

{\Large {\sc BAT} - a short introduction}

\vspace{1.0cm}

{A.~Caldwell$^1$, D.~Koll\'ar$^1$, K.~Kr\"oninger$^2$}

\vspace{1.0cm}

{\it $^1$~Max Planck Institute for Physics, Munich, Germany \\
$^2$~II. Institute of Physics, University of G\"ottingen, G\"ottingen }
\end{center}

\thispagestyle{empty}

\pagebreak

% --------------------------------------------------------
% table of contents
% --------------------------------------------------------

\thispagestyle{empty}

\tableofcontents

\pagebreak

% --------------------------------------------------------
% introduction
% --------------------------------------------------------
\section{Introduction}
\label{section:introduction}

{\sc BAT} is a package for data analysis. It's main purpose is to
handle different models in a Bayesian sense and to help optimize
parameters, and check the validity of the model with the data. An
introduction to {\sc BAT} can be found in~\cite{Caldwell:2008fw}.

The following introduces the basic principals {\sc BAT} is based on
and summarizes how to use the program. This is followed by a brief
example distributed with the software.

% --------------------------------------------------------
% installation
% --------------------------------------------------------
\section{Installation}
\label{section:installation}

See \verb|doc/INSTALL| for installation instructions.

% --------------------------------------------------------
% Running BAT
% --------------------------------------------------------
\section{Running BAT}
\label{section:runningBAT}

This section describes the basic features of {\sc BAT} and the general
flow of work. More complex operations are discussed in a later
section.

\subsection{Getting started}

In {\sc BAT}, all models are represented as individual classes. They
inherit from the same base class which is contained in the {\sc BAT}
library. This library can be linked against in any existing C++
code. In order to start a new project several files need to be
provided by the user:
%
\begin{itemize}
\item A definition of each model used in the analysis, i.e. all source
  files of classes beloning to these models;
\item A main file in which the actual analysis is performed;
\item A makefile which links against the {\sc BAT} library.
\end{itemize}

\noindent
Examples for the these three files come with the distribution of {\sc
BAT} (see section~\ref{section:example}).

\subsection{Creating a model}

All model classes inherit from the abstract base class
\verb|BCModel|. It offers a number of virtual methods which need to be
overloaded by the user. Apart from the constructor and destructor,
these methods are
%
\begin{verbatim}
double BCModel::LogLikelihood(std::vector <double> par)
\end{verbatim}

\noindent
This method contains a mathematical expression for the conditional
probability density (or {\it likelihood}) of the data given a set of
parameter values. For reasons of numerical stability it returns the
logarithm of that expression. \\

\begin{verbatim}
double BCModel::LogAPrioriProbability(std::vector <double> par)
\end{verbatim}
%
\noindent
This method contains a mathematical expression for the {\it a priori}
probability density for a set of parameter values. For reasons of
numerical stability it returns the logarithm of that expression. \\

\noindent
The vector \verb|par| has as many entries as parameters are defined in
the model. The indeces of the elements of the vector correspond to the
indeces of the parameters. The values of these elements are
constrained to the boundaries of the corresponding parameter.

\subsection{Definition of parameters}

The parameters of a model are represented as instances of the class
\verb|BCModelParameter|. The constructor of this class has as
parameters the name of the parameter (\verb|const char * name|), and
the interval for which the parameter is defined (\verb|double xmin|,
\verb|double xmax|). \\ 

Instances of that class can be added to the model via 
%
\begin{verbatim}
void BCModel::AddParameter(BCParameter * parameter) 
\end{verbatim}
%
\noindent 
or they are implicitly defined via 
%
\begin{verbatim}
void BCModel::AddParameter(const char * name, double xmin, double xmax) . 
\end{verbatim}

Every parameter has to have a unique name and valid limits. Once added
to a model, a parameter can be referenced by an index or by its
name. The index is given once the parameter is added to the model,
starting from~0 for the first one. 

\subsection{Data}

In {\sc BAT}, data are described by {\it data points} which are
combined to {\it data sets}. Both are represented by instances of the
classes \verb|BCDataPoint| and \verb|BCDataSet|, respectively.  Data
can be loaded from a plain ASCII files or ROOT files containing a
tree. The following interfaces can be used 
%
\begin{verbatim}
int BCDataSet::ReadDataFromFile(const char* fn, const char* tn, const char* bn)
int BCDataSet::ReadDataFromFile(const char* fn, int n)
\end{verbatim} 
%
\noindent 
where \verb|fn| is the name of the file, \verb|tn| is the name of the
ROOT tree, and \verb|nb| is the name of the branches separated by
commas. The integer \verb|n| defines the number of events which are
read from the ASCII file. The function returns an error code. Once a
data set is defined it can be assigned to a model with
%
\begin{verbatim}
void BCModel::SetDataSet(BCDataSet* dataset) .
\end{verbatim} 

\subsection{Optimization}

In the context of {\sc BAT}, optimization refers to finding the global
mode of the posterior pdf. It can be performed via the member function 
%
\begin{verbatim}
void BCModel::FindMode(); 
\end{verbatim}

Different methods can be used to perform the optimization. They can be
chosen via 
%
\begin{verbatim}
void BCModel::SetOptimizationMethod(BCIntegrate::BCOptimizationMethod method)
\end{verbatim}
%
\noindent
where \verb|BCIntegrate::BCOptimizationMethod| is one of the
following:
%
\begin{itemize}
\item \verb|kOptSimulatedAnnealing| using a simulated annealing
  algorithm,
\item \verb|kOptMetropolis| using the Metropolis algorithm, i.e.,
  optimization is done during the MCMC run, and 
\item \verb|kOptMinuit| using the ROOT version of Minuit. 
\end{itemize}

The default setting ist \verb|kOptMinuit|.

\subsection{Marginalization}

The posterior pdf can be marginalized with respect to a single
parameter or with respect to two parameters. The marginalization can
be performed via: 
%
\begin{verbatim}
int BCModel::MarginalizeAll() . 
\end{verbatim}
%
The function returns the number of marginalized
distributions. Different algorithms can be used, and set via
%
\begin{verbatim}
BCModel::SetMarginalizationMethod(BCIntegrate::BCMarginalizationMethod method), 
\end{verbatim}
%
\noindent 
where \verb|BCIntegrate::BCMarginalizationMethod| is one of the
following:
%
\begin{itemize}
\item \verb|kMargMonteCarlo| using random sampling through the parameter space, and
\item \verb|kMargMetropolis| using Markov Chain Monte Carlo. 
\end{itemize}

The default setting is \verb|kMargMetropolis|. 

\subsection{Goodness-of-fit}

A goodness-of-fit test for the parameters found in the optimization
can be performed using the $p$-value defined
in~\cite{Caldwell:2008fw}. In the most general case, the user has to
supply data sets generated with the parameters to be tested. The
$p$-value can then be evaluated via
%
\begin{verbatim}
BCH1D * GoodnessOfFitTest(const char * filenname, std::vector <double> parameters) ,
\end{verbatim}
%
\noindent 
where the file \verb|filename| contains a list of .txt files which
contain the data sets, and \verb|parameters| is a vector of the
parameters used to generate the data sets. The function returns a
\verb|BCH1D| histogram containing the frequency distribution of the
(logarithm of the) conditional probablity. The $p$-value can be
obtained using 
%
\begin{verbatim}
double BCH1D::GetPValue(double probability) , 
\end{verbatim}
%
\noindent
where probability is the conditional probablity of the pdf (using the
original dat set) at the best fit parameters. 

\subsection{Output: log files}

The class \verb|BCLog| is used to write out information during the
runtime of the program. The output is written to screen and to a log
file. The minimum level of detail can be set independently for both
via
%
\begin{verbatim} 
void BCLog::SetMinimumLogLevelFile(BCLog::LogLevel loglevel) , 
void BCLog::SetMinimumLogLevelScreen(BCLog::LogLevel loglevel) , 
\end{verbatim} 
%
\noindent 
where the level is one of the following 
%
\begin{itemize} 
\item \verb|debug|: Lowest level of information. 
\item \verb|detail|: Details of functions, such as the status of the
  Markov chains, etc.
\item \verb|summary|: Results, such as best-fit values, normalization, etc. 
\item \verb|warning|: Warning messages 
\item \verb|nothing|: Nothing is written out. 
\end{itemize} 

A log file has to be opened in the beginning of the main file using
%
\begin{verbatim} 
void BCLog::OpenLog() . 
\end{verbatim} 

\subsection{Output: summary print-out}

A summary of the results of the Markov Chain Monte Carlo is written to
file at the end of each run. A single file is created for each model
used in the analysis. 

\subsection{Output: histograms and files}

The marginalized distributions of a model can be written to file via
%
\begin{verbatim}
void PrintAllMarginalized(const char * file, int hdiv=1, int ndiv=1) ,
\end{verbatim}
%
\noindent 
where \verb|file| is the filename. 

% --------------------------------------------------------
% Technical tools 
% --------------------------------------------------------
\section{Technical tools}
\label{section:tools}

{\it Save all this for a later version of the manual.} 

\subsection{Optimization}

\subsection{Markov Chain Monte Carlo}

\subsection{Integration}

% --------------------------------------------------------
% More features
% --------------------------------------------------------
\section{More features}
\label{section:features}

{\it Save all this for a later version of the manual.} 

\subsection{Normalization}

\subsection{Model comparison}

\subsection{Fitting}

\subsection{Single event analysis}

\subsection{User defined data interface} 


% --------------------------------------------------------
% Examples 
% --------------------------------------------------------
\section{Example} 
\label{section:example}

Properly commented examples on how to run {\sc BAT} are part of the
distribution. They can be found in \verb|examples/|.

% --------------------------------------------------------
% bibliography 
% --------------------------------------------------------

\addcontentsline{toc}{section}{Bibliography}

\begin{thebibliography}{99}
\bibitem{Caldwell:2008fw}
  A.~Caldwell, D.~Kollar and K.~Kroeninger,
  ``BAT - The Bayesian Analysis Toolkit,''
  arXiv:0808.2552 [physics.data-an].
\end{thebibliography}

% --------------------------------------------------------

\end{document} 


